\documentclass[11pt]{article}
\usepackage[margin=1in]{geometry}
\usepackage{amsmath,amssymb}
\usepackage{booktabs}
\usepackage{graphicx}
\usepackage{hyperref}
\usepackage{microtype}

\title{Physics-Guided Multi-Scale Low-Rank Latent PDE Operator}
\author{Jiawei Zheng}
\date{November 2026}

\begin{document}
\maketitle

\begin{abstract}
We study a physics-guided neural operator for parameterized one-dimensional
partial differential equations. The operator combines a convolutional and
Transformer context encoder, a low-rank latent PDE rollout, multi-scale skip
features, and a grid decoder. Lightweight output-side PDE refiners can be
attached without retraining the backbone. On PDEBench Burgers with
$\nu=0.01$, the Burgers refiner reaches test L2 $0.022853$ and normalized PDE
residual $0.244626$, compared with PINO's $0.024261$ and $0.253304$. A
separate advection output refiner reaches test L2 $0.005789$ and residual
$0.102403$ at $\beta=1$, compared with PINO's $0.006619$ and $0.104517$. An
independent audit shows that the Burgers advantage is confined to boundary and
overall residual: PINO remains better in the time interior and on a
fourth-order residual. The advection advantage persists across overall,
boundary, interior, and high-order residuals. Paired tests over 1,000 common
test samples remain significant after Holm correction. Zero-shot Burgers
transfer to $\nu=0.1$ fails, while the advection refiner transfers to
$\beta\le 0.7$ but not to larger velocities.
\end{abstract}

\section{Introduction}
Neural operators approximate solution maps rather than individual PDE
solutions. Fourier neural operators (FNO) and DeepONet provide strong
parameter-matched baselines \cite{li2021fno,lu2021deeponet}. Physics-informed
variants add residual losses during training \cite{li2021pino}. We investigate
whether a low-rank latent representation plus a small post-hoc physical
correction can improve both data accuracy and physical consistency without
large parameter growth.

\section{Method}
The backbone is
\begin{equation}
u_0(x) \rightarrow \text{context}
\rightarrow z_0
\rightarrow \text{latent PDE rollout}
\rightarrow \text{low-rank decoder}
\rightarrow u_n(x,t).
\end{equation}
For known constant periodic transport, we use an analytic prior
$u_p(x,t)=u_0((x-\beta t)\bmod 1)$.

The Burgers derivative refiner receives
$\big[r/s,\ u_x/\mathrm{rms}(u_x),\ u_{xx}/\mathrm{rms}(u_{xx})\big]$,
where $r=u_t+uu_x-\nu u_{xx}$. The advection output refiner receives the
corresponding $r=u_t+\beta u_x$ features after analytic transport. Both use a
zero-initialized $3\rightarrow h$, $1\times1$, SiLU, $h\rightarrow1$,
$5\times5$ branch.

\section{Experiments}
\begin{table}[h]
\centering
\caption{Primary official-test results. Values are three-seed means.}
\begin{tabular}{llrrr}
\toprule
Task & Model & L2 & residual & parameters \\
\midrule
Burgers & PG + refiner & 0.022853 & 0.244626 & 502,989 \\
Burgers & PINO & 0.024261 & 0.253304 & 502,997 \\
Advection $\beta=1$ & PG + refiner & 0.005789 & 0.102403 & 501,668 \\
Advection $\beta=1$ & PINO & 0.006619 & 0.104517 & 503,165 \\
\bottomrule
\end{tabular}
\end{table}

\paragraph{Residual audit.}
For Burgers, the interior residual is $0.227397$ for PG and $0.221918$ for
PINO; the fourth-order residual is $1.372452$ and $1.367258$, respectively.
For advection $\beta=1$, PG is lower in overall, boundary, interior, and
fourth-order residual. The paired-test Holm-adjusted $p$-values for the four
primary comparisons range from $5.85\times10^{-10}$ to $1.26\times10^{-3}$.

\paragraph{External checks.}
On official Burgers $\nu=0.1$, zero-shot PG and FNO L2 are both about $0.39$.
A 192-parameter viscosity-conditioned joint fine-tune reduces the new
viscosity L2 to $0.0388$, but degrades the original $\nu=0.01$ L2 from $0.0197$
to $0.0264$. The advection $\beta=1$ refiner transfers to $\beta\le0.7$;
larger $\beta$ requires a specialized branch.

\section{Scope and Limitations}
The current evidence supports a bounded claim: approximately 501K parameters,
parameter-matched competitiveness with FNO, and carefully audited overall
metrics against PINO. It does not support universal PDE solving, full NS
momentum accuracy, or unrestricted inverse problems. Diffusion-sorption
parameters were recovered from official metadata, but the current PG residual
($0.999855$) remains above the observed-data residual ($0.531420$).

\section{Conclusion}
Low-rank latent PDE backbones and small output-side physics refiners provide a
controlled way to trade data accuracy against physical residual. The most
robust evidence is the audited advection $\beta=1$ result; Burgers requires a
multi-viscosity training study, and diffusion-sorption requires a
structure-level physical correction.

\bibliographystyle{plain}
\bibliography{references}
\end{document}
